% ID: SEL-D0043-OB-E02
% ESTADO: cotejada. Segunda pasada visual y borrador compilado revisados.
% Fuente: Selectividad/sel_2014_matematicas_aplicadas/EXAMEN -ANDAL-2.pdf
% SHA-256: fc025a25faa338cbb9d8d288443f0322dfbcc0464f7474a16a84e8dfe47f81d1
% Página: 2; opción B; ejercicio original 2.

El porcentaje de personas que sintonizan un programa de radio que se emite entre las 6 y las 12 horas viene dado, según la hora $t$, mediante la función
\[
S(t)=660-231t+27t^2-t^3,\qquad6\leq t\leq12.
\]
\begin{enumerate}[label=\alph*)]
\item \textbf{(0.5 puntos)} \textquestiondown{}Qué porcentaje de personas sintonizan el programa al comenzar la emisión? \textquestiondown{}Y al cierre?
\item \textbf{(2 puntos)} \textquestiondown{}A qué hora tiene máxima y mínima audiencia? \textquestiondown{}Qué porcentaje de personas sintonizan el programa a dichas horas?
\end{enumerate}
